\documentclass[12pt, a4paper]{article}

% ==========================================
% PAQUETES FUNDAMENTALES
% ==========================================
\usepackage[utf8]{inputenc}
\usepackage[T1]{fontenc}
\usepackage[spanish,es-tabla]{babel}
\usepackage{geometry}
\geometry{top=2.5cm, bottom=2.5cm, left=2.5cm, right=2.5cm}
\usepackage{microtype}
\usepackage{tgtermes}
\usepackage[scaled=0.92]{helvet}
\usepackage{graphicx}
\usepackage{tabularx}
\usepackage{booktabs}
\usepackage{caption}
\usepackage{parskip}
\usepackage{enumitem}
\usepackage{setspace}
\usepackage{tcolorbox}
\usepackage{xcolor}
\usepackage{tikz}
\usepackage{titletoc} 
\usepackage{fancyhdr} 
\usepackage{anyfontsize}

% ==========================================
% CONFIGURACIÓN DE COLORES
% ==========================================
\definecolor{unmsm_red}{RGB}{133, 45, 58} % Color vino institucional #852D3A
\definecolor{unmsm_darkred}{RGB}{165, 42, 53} % Granate San Marcos #A52A35
\definecolor{bg_marfil}{RGB}{251, 247, 239} % Fondo marfil oficial #FBF7EF
\definecolor{bg_arena}{RGB}{237, 225, 207} % Fondo arena #EDE1CF
\definecolor{unmsm_gray}{RGB}{106, 98, 90} % Gris calido #6A625A
\definecolor{unmsm_gold}{RGB}{200, 190, 176} % Separadores y acentos #C8BEB0

% ==========================================
% CONFIGURACIÓN DE TÍTULOS (Estilo visual)
% ==========================================
\usepackage[explicit]{titlesec}
\usepackage{hyperref}

% Formato para CAPÍTULO
\titleformat{\section}[display]
  {\normalfont\huge\bfseries\color{black}}
  {\vspace*{0.1cm} % <-- AÑADE ESTA LÍNEA AQUÍ PARA BAJAR EL TÍTULO
   \colorbox{unmsm_red}{\textcolor{white}{\textsf{\textbf{CAPÍTULO}}}} \hspace{0.2cm} \textcolor{unmsm_red}{\fontsize{60}{60}\selectfont\thesection}}
  {1ex}
  {\Huge #1}

% Formato para Subtemas 1.1
\titleformat{\subsection}
  {\normalfont\Large\sffamily\bfseries\color{black}}
  {\thesubsection}
  {0.5em}
  {#1}
  [\vspace{0.1cm}\textcolor{black}{\rule{\textwidth}{0.8pt}}]

% Formato para Apartados específicos y Casos 1.1.1
\titleformat{\subsubsection}
  {\normalfont\large\sffamily\bfseries\color{black}}
  {}
  {0em}
  {#1}
  [\vspace{0.1cm}\textcolor{black}{\rule{\textwidth}{0.4pt}}]

% ==========================================
% COMANDOS PERSONALIZADOS (Cajas, Notas y Listas)
% ==========================================
% Espaciado global de listas para optimizar maquetacion
\setlist[itemize]{topsep=0.3em, itemsep=0.25em, parsep=0pt, leftmargin=1.5em}
\setlist[enumerate]{topsep=0.3em, itemsep=0.25em, parsep=0pt, leftmargin=1.5em}

\newtcolorbox{cajabeige}{
    colback=bg_arena, colframe=bg_arena, arc=3pt, boxrule=0pt,
    top=0.4cm, bottom=0.4cm, left=0.6cm, right=0.6cm
}

\newcommand{\textoside}[2]{
    \vspace{0.25cm}
    \noindent
    \begin{minipage}[t]{0.68\textwidth}
        #1
    \end{minipage}%
    \hfill%
    \begin{minipage}[t]{0.28\textwidth}
        \begin{tcolorbox}[colback=bg_arena, colframe=unmsm_gold, arc=3pt, boxrule=0.6pt, top=0.3cm, bottom=0.3cm, left=0.35cm, right=0.35cm]
            \raggedright
            \textcolor{unmsm_red}{\scriptsize\sffamily\textbf{#2}}
        \end{tcolorbox}
    \end{minipage}
    \vspace{0.35cm}
}

\newlist{resumenlist}{enumerate}{1}
\setlist[resumenlist]{
    label=\textcolor{unmsm_red}{\Huge 0\arabic*},
    leftmargin=1.3cm, labelsep=0.6cm, itemsep=0.4cm, topsep=0.3cm
}

% ==========================================
% CONFIGURACIÓN DEL ÍNDICE
% ==========================================
\titlecontents{section}[0em]
  {\vspace{0.35cm}\large\bfseries\sffamily\color{unmsm_red}}
  {\thecontentslabel\hspace{0.7em}}
  {}
  {\titlerule*[0.6pc]{.}\textbf{\contentspage}}
  [\vspace{0.1cm}]

\titlecontents{subsection}[1.8em]
  {\vspace{0.12cm}\normalsize\sffamily\color{black}}
  {\thecontentslabel\hspace{0.6em}}
  {}
  {\titlerule*[0.6pc]{.}\contentspage}
  []

% ==========================================
% CONFIGURACIÓN DE ENCABEZADOS Y PIES (fancyhdr)
% ==========================================
\pagestyle{fancy}
\fancyhf{} 
% Encabezado (Header institucional estilo San Marcos)
\fancyhead[L]{\textcolor{unmsm_red}{\small\sffamily\textbf{\leftmark}}} 
\fancyhead[R]{\textcolor{unmsm_red}{\small\sffamily Creaci\'on de Prototipos R\'apidos \quad\rule[-0.2ex]{0.5pt}{1.2ex}\quad \textbf{\thepage}}} 
\renewcommand{\headrule}{\hbox to\headwidth{\color{unmsm_red}\leaders\hrule height 0.6pt\hfill}}
% Pie de pagina (Footer limpio, sin duplicacion de numero de pagina)
\fancyfoot[L]{\textcolor{unmsm_gray}{\scriptsize\sffamily Universidad Nacional Mayor de San Marcos}} 
\fancyfoot[R]{\textcolor{unmsm_gray}{\scriptsize\sffamily Inteligencia Artificial Aplicada}} 
\renewcommand{\footrulewidth}{0pt} 

% Sobrescribir el estilo 'plain'
\fancypagestyle{plain}{
    \fancyhf{}
    \fancyhead[L]{\textcolor{unmsm_red}{\small\sffamily\textbf{\leftmark}}}
    \fancyhead[R]{\textcolor{unmsm_red}{\small\sffamily Creaci\'on de Prototipos R\'apidos \quad\rule[-0.2ex]{0.5pt}{1.2ex}\quad \textbf{\thepage}}}
    \renewcommand{\headrule}{\hbox to\headwidth{\color{unmsm_red}\leaders\hrule height 0.6pt\hfill}}
    \fancyfoot[L]{\textcolor{unmsm_gray}{\scriptsize\sffamily Universidad Nacional Mayor de San Marcos}}
    \fancyfoot[R]{\textcolor{unmsm_gray}{\scriptsize\sffamily Inteligencia Artificial Aplicada}}
    \renewcommand{\footrulewidth}{0pt}
}

\hypersetup{colorlinks=true, linkcolor=unmsm_red, urlcolor=cyan, pdftitle={Creación de Prototipos Rápidos}}

\begin{document}

% ==========================================
% INICIO DE NUMERACIÓN ROMANA (Preliminares)
% ==========================================
\pagenumbering{roman} % Páginas i, ii, iii...
\pagecolor{bg_marfil} % Fondo marfil global

% ==========================================
% PORTADA (Página i - Oculta)
% ==========================================
\begin{titlepage}
    \color{unmsm_red} 
    \centering
    
    \vspace*{-1cm}
    \begin{tikzpicture}
        \node[inner sep=0pt] at (0,0) {\includegraphics[width=6.0cm]{escudo-unmsm.png}}; 
    \end{tikzpicture}\\[0.8cm]
    
    {\LARGE \textbf{UNIVERSIDAD NACIONAL}}\\[0.15cm]
    {\LARGE \textbf{MAYOR DE SAN MARCOS}}\\[0.6cm]
    {\large Facultad de Ciencias Matemáticas}\\[0.4cm]
    \noindent\rule{5cm}{0.4pt}\\[0.6cm] 
    {\large INTELIGENCIA ARTIFICIAL Y COMPUTACIÓN CIENTÍFICA}\\[1.5cm]
    
    {\Huge \textbf{Creación de Prototipos Rápidos:}}\\[0.4cm]
    {\LARGE \textbf{Guía de Herramientas No-Code/Low-Code}}\\[0.2cm]
    {\LARGE \textbf{para Validar Ideas Científicas}}\\[1.2cm]
    
    \parbox{0.85\textwidth}{\centering \large \textit{Marcos Metodológicos, Arquitecturas de Microservicios y Gobernanza con Inteligencia Artificial Generativa}}\\[1cm]
    
    \vfill
    \begin{tikzpicture}[remember picture,overlay]
        \fill[bg_arena] (current page.south west) rectangle ([yshift=5.5cm]current page.south east);
        \node[align=center, text=unmsm_red] at ([yshift=3.5cm]current page.south) {
            {\normalsize Grupo de Investigación en Inteligencia Artificial y Computación Científica, UNMSM}\\[0.3cm]
            {\normalsize Asignatura de Inteligencia Artificial Aplicada}
        };
        \node[align=center, text=unmsm_red] at ([yshift=1.2cm]current page.south) {
            {\small LIMA, PERÚ $\cdot$ 2026}
        };
    \end{tikzpicture}
\end{titlepage}

\color{black} % Restaurar color de texto general

% ==========================================
% DATOS DE LA OBRA (Página ii)
% ==========================================
\newpage
\markboth{DATOS DE LA OBRA}{DATOS DE LA OBRA} 
\phantomsection
\addcontentsline{toc}{section}{Datos de la Obra}

\vspace*{0.5cm}
{\LARGE \textbf{Creación de Prototipos Rápidos: Guía de Herramientas No-Code/Low-Code para Validar Ideas Científicas}}\\[0.5cm]
{\large \textcolor{darkgray}{Marcos Metodológicos, Arquitecturas de Microservicios y Gobernanza con Inteligencia Artificial Generativa}}

\vspace{1.2cm}
{\Large \textsf{\textbf{Autoría}}}\\[-0.2cm]
\textcolor{gray}{\rule{\textwidth}{0.8pt}}\\[0.3cm]
{\normalsize Pablo Alexis Valdivia Alipazaga}\\[0.1cm]
{\normalsize Pablo Gerardo Valdivia Alipazaga}\\[0.1cm]
{\normalsize Carlos Arnaldo Huamani Riojas}\\[0.1cm]
{\normalsize Adriano Collantes Sánchez}\\[0.1cm]
{\normalsize Guillermo Valentino Ayora Príncipe}

\vspace{1.2cm}
{\Large \textsf{\textbf{Información académica}}}\\[-0.2cm]
\textcolor{gray}{\rule{\textwidth}{0.8pt}}\\[0.3cm]
{\normalsize Asignatura de Inteligencia Artificial Aplicada}\\[0.1cm]
{\normalsize Docente: Frida Mereyda Lopez Cordova}\\[0.1cm]
{\normalsize Ciclo Académico 2026-II}\\[0.1cm]
{\normalsize Lima, Perú $\cdot$ 2026}

\vfill
\noindent
{\small \textcolor{darkgray}{Documento académico elaborado como parte de la evaluación de la asignatura de Inteligencia Artificial Aplicada.}}\\[0.2cm]
{\small \textcolor{darkgray}{Derechos de autor \copyright~2026. Se permite su uso y consulta para fines estrictamente académicos e investigativos dentro de la Universidad Nacional Mayor de San Marcos.}}
\vspace{1cm}

% ==========================================
% CONTENIDOS / ÍNDICE (Página iii)
% ==========================================
\newpage
\markboth{CONTENIDOS}{CONTENIDOS} 
\vspace*{0.5cm}
{\Huge \textbf{Contenidos}}\\[0.8cm]

\makeatletter
\@starttoc{toc}
\makeatother


% ==========================================
% INICIO DE NUMERACIÓN ARÁBIGA (Cuerpo Principal)
% ==========================================
\newpage
\pagenumbering{arabic} % Reinicia en 1 y usa números arábigos

% ==========================================
% PRESENTACIÓN (Página 1)
% ==========================================
\markboth{PRESENTACIÓN}{PRESENTACIÓN}
\phantomsection
\addcontentsline{toc}{section}{Presentación}

\vspace*{0.5cm}
{\Huge \textbf{Presentación}}
\vspace{0.8cm}

La validación empírica de hipótesis científicas mediante software ha enfrentado históricamente cuellos de botella como la escasez de personal especializado y prolongados ciclos de codificación. La convergencia de las plataformas de bajo código (Low-Code/No-Code, LCNC) y la Inteligencia Artificial Generativa (IAG) impulsa el \textbf{prototipado científico acelerado}, un paradigma centrado en la reasignación estratégica de la carga cognitiva (\textit{task reallocation}). Bajo este enfoque, la automatización asume la construcción sintáctica repetitiva mientras el investigador reconcentra su esfuerzo en el pensamiento sistémico, la arquitectura y la verificación de calidad.

\vspace{0.4cm}
Esta obra, estructurada bajo la plantilla formal de la Universidad Nacional Mayor de San Marcos (UNMSM), ofrece una guía rigurosa fundamentada en la evidencia empírica de cinco investigaciones internacionales clave (2023–2026). A lo largo de sus cinco capítulos, se aborda el diseño guiado por modelos (MDD), la automatización de microservicios desde diagramas UML (FSMicroGenerator), la gobernanza documental bajo normas ISO 15489/27001, el \textit{Vibe Coding} respaldado por el marco VIBE4M y la orquestación de la autonomía mediante la Escalera de Capacidades (\textit{Capability Ladder}).

\vspace{0.4cm}
Dirigido a docentes, investigadores, ingenieros y estudiantes de la UNMSM y de la comunidad académica internacional, este libro constituye un manual de referencia metodológica para liderar la innovación tecnológica con rigor científico y responsabilidad ética.

\vfill

\begin{flushright}
    \textcolor{darkgray}{Grupo de Investigación en Inteligencia Artificial y Computación Científica, UNMSM}\\
    \textcolor{darkgray}{Lima, Perú $\cdot$ 2026}
\end{flushright}
\vspace*{0.5cm}

% ==========================================
% CAPÍTULO 1
% ==========================================
\newpage
\section[Fundamentos y Gobernanza Documental en LCNC]{Fundamentos, Abstracción Visual y Gobernanza Documental en Prototipado Científico}
\markboth{CAPÍTULO 1\quad FUNDAMENTOS Y GOBERNANZA}{}

\begin{cajabeige}
\textbf{Pregunta Central del Capítulo:} ¿Cómo transforman las plataformas Low-Code/No-Code los paradigmas tradicionales de desarrollo de software para acelerar el prototipado científico, garantizando al mismo tiempo el cumplimiento de estándares internacionales de gobernanza documental y control de acceso?\\[0.3cm]
\textbf{Objetivo de Aprendizaje:} El lector comprenderá la evolución de las metodologías de desarrollo, identificará los principios de abstracción visual y modelado guiado por modelos (MDD), y evaluará la arquitectura, eficiencia de costos y cumplimiento normativo de las plataformas de bajo código en comparación con arquitecturas tradicionales basadas en código.
\end{cajabeige}

\textoside{
La democratización de la ingeniería de software es el proceso mediante el cual se reducen las barreras técnicas de entrada en la creación de aplicaciones mediante interfaces gráficas, componentes reutilizables y automatización declarativa. Este paradigma permite a los desarrolladores ciudadanos (\textit{citizen developers}) —investigadores, científicos y analistas sin formación profunda en sintaxis de programación— construir, probar y desplegar prototipos funcionales robustos sin depender exclusivamente de equipos centrales de TI.
}{CONCEPTO CLAVE\\[0.1cm]Democratización de la Ingeniería de Software}
\vspace{0.5cm}

\subsection{Evolución Metodológica del Desarrollo de Software y Abstracción Visual}

La historia de la ingeniería de software está marcada por una búsqueda continua de niveles superiores de abstracción. Como explican Sido y Emon (2024), las metodologías tradicionales orientadas a procesos rígidos —como el Modelo en Cascada (\textit{Waterfall}), el Modelo en Espiral de Boehm y el Modelo en V— estaban diseñadas para entornos con requerimientos estáticos y predecibles. Aunque estos enfoques tradicionales proporcionaban un control de gestión detallado y una documentación exhaustiva, generaban ciclos de desarrollo sumamente estructurados.

La transición hacia metodologías ágiles (como Scrum y Extreme Programming) y el Proceso Unificado Rational (RUP) mitigó parcialmente estos cuellos de botella al introducir la iteración incremental y la adaptación continua. Sin embargo, como señalan Khalfaoui, Khalfaoui y Azmani (2025), el desarrollo de aplicaciones web full-stack modernas sigue enfrentando una sobrecarga cognitiva considerable. Estudios de productividad en ingeniería de software revelan que los desarrolladores dedican aproximadamente el 80\% de su tiempo a tareas de depuración y corrección de errores, de los cuales entre el 40\% y el 50\% corresponde a retrabajo evitable derivado de la codificación manual repetitiva de sintaxis \textit{boilerplate}.

Para superar la sobrecarga de la codificación textual directa, surgieron las plataformas de desarrollo de bajo código y sin código (Low-Code/No-Code Development Platforms, LCDPs/NCDPs). Estas plataformas operan bajo el principio del Diseño Guiado por Modelos (Model-Driven Design, MDD), en el cual la arquitectura, la lógica de negocio y las interfaces de usuario se definen mediante modelos conceptuales gráficos de alto nivel. El motor de la plataforma traduce automáticamente estas especificaciones visuales en código fuente ejecutable o configuraciones interpretadas en tiempo de ejecución.

Como destacan Reudismam Rolim et al. (2026), el paradigma low-code democratiza el desarrollo de software al proporcionar interfaces visuales, componentes prediseñados y abstracciones de alto nivel que permiten a los "desarrolladores ciudadanos" e investigadores participar activamente en la construcción de aplicaciones funcionales sin requerir una formación técnica profunda en sintaxis de programación.

Sido y Emon (2024) destacan que la distinción fundamental entre el desarrollo Low-Code y No-Code radica en la audiencia objetivo y el nivel de extensibilidad técnica:

\begin{itemize}
    \item \textbf{Plataformas No-Code:} Orientadas a usuarios de negocio y científicos sin conocimientos de programación. Utilizan interfaces puramente visuales de arrastrar y soltar (\textit{drag-and-drop}) con bloques funcionales cerrados. Son ideales para prototipos de complejidad baja a moderada, formularios de recolección de datos y automatizaciones de flujos simples.
    \item \textbf{Plataformas Low-Code:} Diseñadas para coordinar el trabajo entre desarrolladores ciudadanos y profesionales informáticos. Aunque ofrecen abstracción visual para el 80\% de la aplicación, permiten la inyección de código personalizado (JavaScript, Python, C\#, SQL) para algoritmos complejos e integraciones personalizadas.
\end{itemize}

\subsection{Arquitecturas de Automatización Empresarial y Gobernanza Documental}

En el ámbito científico y académico, el prototipado rápido no solo requiere la generación ágil de interfaces, sino también un estricto control de la gobernanza documental, la trazabilidad de decisiones y la auditoría de procesos. La investigación desarrollada por Marian, Neferu y Mitrea (2026) demuestra cómo las plataformas Low-Code empresariales integradas en la nube —específicamente el ecosistema Microsoft Power Platform— proporcionan un entorno robusto para operacionalizar flujos de trabajo documentales complejos bajo estándares de cumplimiento industrial.

A diferencia de los entornos de codificación tradicional y del \textit{vibe coding} asistido puramente por inteligencia artificial estocástica, las plataformas low-code ofrecen un marco estructurado que equilibra accesibilidad y rigurosidad técnica. El análisis sistemático formulado por Reudismam Rolim et al. (2026) evidencia que las plataformas de bajo código mantienen un control arquitectónico estructurado, dado que el diseño de las soluciones se encuentra delimitado y gobernado por la arquitectura y los esquemas deterministas de la plataforma proveedora.

Asimismo, Reudismam Rolim et al. (2026) subrayan que el desarrollo low-code reduce sustancialmente los riesgos de seguridad en comparación con el código generado de manera probabilística o manual sin auditorías exhaustivas. Al heredar protocolos nativos de gobernanza, autenticación empresarial y capas de validación preconfiguradas, se minimiza la superficie de vulnerabilidades ocultas y la acumulación acelerada de deuda técnica en los entornos experimentales de I+D.

La arquitectura de una solución Low-Code empresarial se organiza conceptualmente en tres capas interconectadas:

\begin{enumerate}
    \item \textbf{Capa de Interfaz de Usuario (UI):} Desarrollada mediante aplicaciones de lienzo o basadas en modelos (Power Apps), ofreciendo formularios interactivos con validación de datos en tiempo real y paneles de control (dashboards) adaptados al rol del usuario.
    \item \textbf{Capa de Lógica de Negocio y Orquestación:} Ejecutada por motores de flujos de trabajo (Power Automate), encargados de gestionar los estados de aprobación, las transiciones departamentales, el envío automatizado de notificaciones y la generación sintética de documentos finales.
    \item \textbf{Capa de Almacenamiento de Datos y Persistencia:} Diseñada mediante repositorios estructurados (SharePoint Online / Microsoft Dataverse), respaldada por autenticación centralizada y control de acceso basado en roles (RBAC) mediante Microsoft Entra ID.
\end{enumerate}

\subsubsection{El Ecosistema Power Platform y las Normas de Gobernanza Documental}

Un aspecto crítico resaltado por Marian, Neferu y Mitrea (2026) es la capacidad de las soluciones de bajo código para garantizar el cumplimiento de normas internacionales de gestión de registros e información. A diferencia de los prototipos construidos \textit{ad hoc} con scripts individuales de Python o PHP, las soluciones Low-Code integradas en ecosistemas de nube empresarial heredan automáticamente controles de gobernanza como:

\begin{itemize}
    \item \textbf{Sellado de Tiempo Digital e Identificadores Únicos:} Cada decisión de aprobación o rechazo en el flujo de trabajo queda registrada de forma inmutable con una marca de tiempo precisa y un identificador de auditoría.
    \item \textbf{Control de Versiones Automático:} Las modificaciones e iteraciones sobre documentos científicos se gestionan de forma nativa en bibliotecas con control de versiones, impidiendo la pérdida de historial de revisión.
    \item \textbf{Control de Acceso Centralizado:} La autenticación exige verificación multifactorial y aplica directivas organizacionales para clasificar la sensibilidad de los datos e impedir la pérdida de activos de investigación.
\end{itemize}

\subsection{Caso o Aplicación: Automatización del Flujo Documental MAF en la Gestión Académico-Empresarial}

Para ilustrar la aplicación práctica de estos conceptos, analicemos el caso de la automatización del proceso de aprobación de solicitudes documentales complejas (\textit{Multiple Attachments Flow}, MAF), documentado por Marian, Neferu y Mitrea (2026). En entornos universitarios y centros de investigación, la aprobación de protocolos de investigación, solicitudes de subvención (\textit{grants}) y autorizaciones de compras de equipamiento científico suele sufrir retrasos debido al intercambio ineficiente de correos electrónicos y formularios impresos.

\subsubsection*{Descripción del Caso y Resultados Empíricos}

El sistema desarrollado integra Power Apps como interfaz de captura estructurada, SharePoint Online como base de datos ligera utilizando columnas de búsqueda (\textit{Lookup Columns}) para simular relaciones relacionales sin el costo de licencias de bases de datos dedicadas, y Power Automate para la orquestación.

\begin{figure}[htpb!]
\centering
\begin{tcolorbox}[colback=white, colframe=unmsm_gold, arc=3pt, boxrule=0.8pt, top=0.3cm, bottom=0.3cm, left=0.5cm, right=0.5cm]
\centering
\small\sffamily
\textbf{Figura 1.1.} Arquitectura del Flujo Documental Multitarea (MAF) en Power Platform\[0.2cm]
\begin{tikzpicture}[node distance=1.2cm, auto, >=stealth, font=\footnotesize\sffamily]
    \node [rectangle, draw=unmsm_red, fill=bg_marfil, rounded corners=2pt, text width=2.6cm, text centered] (papps) {\textbf{Power Apps}\\Captura UI / RBAC};
    \node [rectangle, draw=unmsm_red, fill=bg_marfil, rounded corners=2pt, text width=2.6cm, text centered, right of=papps, node distance=3.8cm] (pauto) {\textbf{Power Automate}\\Orquestador MAF};
    \node [rectangle, draw=unmsm_red, fill=bg_marfil, rounded corners=2pt, text width=2.6cm, text centered, right of=pauto, node distance=3.8cm] (sp) {\textbf{SharePoint / Dataverse}\\Persistencia \& ISO};
    \node [rectangle, draw=unmsm_darkred, fill=bg_arena, rounded corners=2pt, text width=2.6cm, text centered, below of=pauto, node distance=1.6cm] (notif) {\textbf{Aprobación Digital}\\Notificación y Firma};
    \draw [->, thick, draw=unmsm_red] (papps) -- (pauto);
    \draw [->, thick, draw=unmsm_red] (pauto) -- (sp);
    \draw [<->, thick, draw=unmsm_darkred] (pauto) -- (notif);
\end{tikzpicture}\\[0.25cm]
\raggedright
\scriptsize \textbf{Fuente:} Adaptación metodológica basada en Marian, Neferu y Mitrea (2026).
\end{tcolorbox}
\vspace{-0.2cm}
\end{figure}

En un estudio piloto controlado de seis semanas con 25 usuarios activos en un entorno académico-industrial, se obtuvieron los siguientes resultados:

\begin{itemize}
    \item \textbf{Reducción del Tiempo de Ciclo:} El tiempo medio de aprobación de solicitudes documentales se redujo en un 65.6\%, pasando de 3.2 días en el proceso manual tradicional a solo 1.1 días con la solución Low-Code automatizada.
    \item \textbf{Usabilidad y Adopción:} El 92\% de los participantes calificó la interfaz de la aplicación como "intuitiva" o "muy intuitiva", y el 88\% completó su primera presentación documental sin requerir asistencia técnica.
\end{itemize}

\subsection{Demostración y Evaluación Comparativa: Power Platform frente a Scripts Tradicionales en Python}

Para comprobar empíricamente las capacidades de automatización e iteración de datos, se contrastó el desempeño de una arquitectura low-code basada en Microsoft Power Platform (Power Apps y Power Automate) frente a la implementación tradicional mediante scripts manuales en Python (utilizando librerías como ReportLab, Pandas y smtplib).

Los resultados empíricos derivados de una evaluación controlada con 25 participantes revelaron las siguientes ventajas determinantes a favor de la plataforma de bajo código:

\begin{itemize}
    \item \textbf{Reducción del tiempo de desarrollo e implementación:} La construcción y puesta en marcha del flujo automatizado en Power Automate demandó un promedio de $\sim$2 horas, en marcado contraste con las $\sim$6 horas requeridas para codificar, depurar y configurar el script en Python.
    \item \textbf{Disminución de la carga de pruebas y depuración:} Power Platform redujo el tiempo total de desarrollo y pruebas en un 80\% a 90\% en comparación con el enfoque tradicional. Mientras que la depuración en Python exigió el rastreo manual de excepciones y líneas de código, Power Automate proporcionó historiales de ejecución detallados por acción y control granular de fallos mediante bloques condicionales visuales (\textit{Configure Run After}).
    \item \textbf{Superioridad en iteración rápida de flujos de datos:} Power Platform superó con holgura a los scripts tradicionales en Python al simplificar la iteración rápida sobre flujos de trabajo de datos, permitiendo reconfigurar reglas lógicas y estructuras documentales desde la interfaz gráfica sin necesidad de refactorizar ni redesplegar código fuente.
    \item \textbf{Seguridad y gobernanza integrada:} El flujo en Power Platform gestionó las credenciales y el cifrado de comunicaciones directamente bajo Azure Active Directory / Microsoft Entra ID y directivas de cumplimiento, superando la gestión manual de credenciales inherente a los scripts de automatización convencionales.
\end{itemize}

% --- NOTAS CAP 1 ---
\newpage
\markboth{CAPÍTULO 1\quad NOTAS BIBLIOGRÁFICAS E HISTÓRICAS}{}
\vspace*{0.8cm}
\noindent{\Huge \textbf{Notas bibliográficas\\[0.2cm] e históricas}}\\[1cm]

El desarrollo del bajo código y la automatización visual hunde sus raíces teóricas en las investigaciones sobre Ingeniería de Software Dirigida por Modelos (MDSE) compiladas por Brambilla, Cabot y Wimmer (2017), así como en el trabajo pionero de James Martin (1991) sobre el Desarrollo Rápido de Aplicaciones (RAD). El desarrollo de este capítulo se fundamenta en las investigaciones de Reudismam Rolim et al. (2026) sobre el análisis sistemático y la comparativa técnica entre desarrollo manual y low-code; en la evaluación empírica de arquitecturas documentales low-code realizada por Marian et al. (2026); y en el estudio de adopción ágil e inteligencia artificial generativa de la Universidad de Aalborg (2024).

% --- RESUMEN CAP 1 ---
\newpage
\markboth{CAPÍTULO 1\quad RESUMEN}{}
\vspace*{0.8cm}
{\Huge \textbf{Resumen del capítulo 1}}
\vspace{1cm}

\begin{resumenlist}
    \item \textbf{Democratización con rigor:} El bajo código traslada el desarrollo hacia los usuarios de dominio y científicos, manteniendo un control arquitectónico estricto.
    \item \textbf{Agilidad de Prototipado:} El uso de herramientas de bajo código permite iterar sobre conceptos científicos y desplegar prototipos funcionales en cuestión de horas.
    \item \textbf{Seguridad por diseño:} Al apoyarse en la infraestructura de plataformas consolidadas, se reducen las vulnerabilidades críticas de seguridad comunes en la programación manual no estandarizada.
    \item \textbf{Eficiencia empírica comprobada:} Power Platform demostró una reducción de tiempos de desarrollo del 60\% al 80\% frente al desarrollo por scripts en Python.
\end{resumenlist}

% --- CONCLUSIONES CAP 1 ---
\newpage
\markboth{CAPÍTULO 1\quad CONCLUSIONES}{}
\vspace*{0.8cm}
{\Huge \textbf{Conclusiones del capítulo 1}}
\vspace{1cm}

\subsubsection*{Conclusión principal}
Las plataformas Low-Code constituyen un catalizador fundamental para la modernización de los procesos de investigación y validación tecnológica. Basado en las evidencias del estudio de la Universidad de Aalborg (2024), el modelado ágil e iterativo intrínseco a las plataformas low-code permite acortar radicalmente los ciclos de desarrollo y optimizar los tiempos de respuesta, facilitando que los investigadores lleven a cabo un prototipado rápido y logren validar hipótesis científicas tempranamente ante comités evaluadores, pares académicos e inversores.

\subsubsection*{Aportes y límites}
\textbf{Aportes:} Se demuestra cuantitativa y conceptualmente cómo la abstracción de componentes acelera la entrega de prototipos funcionales y preserva la trazabilidad de los datos sin comprometer la seguridad. \textbf{Límites:} Las plataformas propietarias dependen de esquemas y modelos de licenciamiento que deben ser ponderados frente a requerimientos analíticos de cómputo intensivo.

\subsubsection*{Preguntas abiertas}
¿Hasta qué punto la integración sinérgica de asistentes de Inteligencia Artificial Generativa dentro de estos entornos low-code permitirá delegar la formulación de consultas complejas sin erosionar el control analítico del investigador?
% ==========================================
% CAPÍTULO 2
% ==========================================
\newpage
\section[Microservicios Automatizados desde UML]{Arquitectura de Microservicios Automatizada desde Modelos UML (FSMicroGenerator)}
\markboth{CAPÍTULO 2\quad MICROSERVICIOS DESDE UML}{}

\begin{cajabeige}
\textbf{Pregunta Central del Capítulo:} ¿De qué manera la generación automatizada de código basada en modelos UML (\textit{FSMicroGenerator}) permite superar los cuellos de botella de escalabilidad, complejidad arquitectónica y mantenimiento en prototipos científicos frente al enfoque monolítico tradicional? \\[0.3cm] %[cite: 1, 2]
\textbf{Objetivo de Aprendizaje:} Comprender la anatomía funcional, técnica y operativa del marco \textit{FSMicroGenerator}, analizar las evidencias cuantitativas/cualitativas de su evaluación experimental y examinar su caso de aplicación real en plataformas de gestión científica e industrial. %
\end{cajabeige}

\textoside{
La arquitectura monolítica agrupa la interfaz de usuario, la lógica de negocio y el acceso a datos en una única unidad desplegable. %[cite: 2] En prototipos científicos de rápido crecimiento, esta rigidez genera un acoplamiento extremo, dificulta el escalado independiente de componentes críticos y provoca que hasta un 80\% del tiempo de desarrollo se consuma en depuración y corrección de errores redundantes. %[cite: 2, 6]
}{CONCEPTO CLAVE\\[0.1cm]Deuda Técnica y Cuellos Monolíticos}
\vspace{0.5cm}

\subsection{Desafíos de Escalabilidad en Prototipos Científicos y el Problema del Monolito}

El desarrollo acelerado de aplicaciones web modernas para la investigación científica enfrenta barreras severas debido al incremento en la complejidad arquitectónica y la necesidad de reducir el tiempo de salida al mercado (\textit{time-to-market}). %[cite: 1] Estudios del \textit{National Institute of Standards and Technology} (NIST) señalan que los desarrolladores dedican aproximadamente el 80\% de su tiempo a la depuración y corrección de código, de los cuales entre el 40\% y el 50\% corresponde a retrabajos totalmente evitables derivados de la codificación manual de tareas repetitivas. %[cite: 6]

En el contexto de la validación de hipótesis científicas, la mayoría de prototipos iniciales se construyen mediante \textbf{arquitecturas monolíticas}. %[cite: 2, 7] Si bien los monolitos facilitan la fase inicial de desarrollo, desencadenan serios problemas a medida que el sistema escala o incorpora procesamiento intensivo de datos: %[cite: 2]

\begin{enumerate}
    \item \textbf{Falta de aislamiento de fallos:} Un error en un módulo secundario puede comprometer la disponibilidad de todo el sistema distribuido. %[cite: 2]
    \item \textbf{Imposibilidad de escalado granular:} No es posible asignar recursos computacionales exclusivamente a los módulos con mayor demanda de procesamiento sin duplicar la aplicación completa. %[cite: 2]
    \item \textbf{Alta barrera de entrada para microservicios:} Aunque la arquitectura de microservicios soluciona estos problemas al dividir la aplicación en servicios independientes y desacoplados, introduce una alta complejidad técnica en orquestación, comunicación interservicio y configuración de infraestructura, sobrecargando a equipos de investigación con personal limitado. %[cite: 1, 2]
\end{enumerate}

\subsection{Anatomía del Marco FSMicroGenerator}

\textoside{
\textit{FSMicroGenerator} aplica los principios de \textit{Domain-Driven Design} (DDD) para descomponer sistemas complejos en contextos delimitados (\textit{bounded contexts}) mediante una jerarquía de tres niveles: \textbf{Paquete} (unidad funcional básica), \textbf{Módulo} (colaboración de paquetes) y \textbf{Proyecto} (sistema distribuido listo para producción). %[cite: 8, 9]
}{CONCEPTO CLAVE\\[0.1cm]Descomposición por Dominios (DDD)}
\vspace{0.5cm}

\subsubsection{Enfoque de Plantillas y Pila Tecnológica}
El motor de generación opera mediante la articulación de dos tipos de plantillas: %[cite: 10]
\begin{itemize}
    \item \textbf{Plantillas Estáticas:} Definen el esqueleto arquitectónico estándar del proyecto, incluyendo la configuración del \textit{API Gateway}, utilidades compartidas, estructura base del \textit{frontend} y la integración de microservicios transversales (gestión de usuarios, autenticación y taxonomía). %[cite: 10, 11]
    \item \textbf{Plantillas Dinámicas:} Transforman los diagramas de clases UML definidos por el usuario en artefactos específicos: esquemas Mongoose, controladores, repositorios, rutas Express.js, pruebas unitarias/integración con Jest, documentación Swagger y componentes \textit{frontend} dinámicos en Angular/PrimeNG. %[cite: 10]
\end{itemize}

\subsubsection{Pila Tecnológica y Ecosistema DevSecOps}
\textit{FSMicroGenerator} unifica el desarrollo mediante la pila \textbf{MEAN} (MongoDB, Express.js, Angular, Node.js), eliminando la fragmentación de lenguajes entre \textit{frontend} y \textit{backend}. %[cite: 15, 16] La infraestructura técnica integra:
\begin{itemize}
    \item \textbf{Comunicación Interservicio:} \textit{RabbitMQ} para mensajería asíncrona y llamadas RPC síncronas de baja latencia. %[cite: 17]
    \item \textbf{Caché y Rendimiento:} \textit{Redis} como capa de almacenamiento temporal para reducir la latencia de lectura. %[cite: 18]
    \item \textbf{Observabilidad Avanzada:} Instrumentación nativa con \textit{Prometheus} (exposición de endpoints \texttt{/metrics}) y tableros interactivos en \textit{Grafana}. %
    \item \textbf{DevSecOps Shift-Left:} Integración nativa en pipelines de \textit{GitLab CI/CD} con análisis estático de código (SAST), escaneo de dependencias, detección de secretos, escaneo de contenedores con \textit{Trivy} e inspección de infraestructura como código (IaC) con \textit{Checkov}. %
\end{itemize}

\subsubsection{Ciclo de Vida del Paquete}
El proceso de construcción consta de cuatro etapas principales: %
\begin{enumerate}
    \item \textbf{Fase de Creación:} Definición semimanual del modelo conceptual mediante diagramas de clases UML (entidades, atributos, relaciones y restricciones) y vista previa en tiempo real. %
    \item \textbf{Fase de Validación:} Revisión técnica por un usuario calificado para detectar incoherencias en el modelo antes de la ejecución. %
    \item \textbf{Fase de Adaptación:} El motor extrae la configuración, despliega las plantillas estáticas y dinámicas, genera el código fuente y realiza un \textit{push} automático a GitLab. %
    \item \textbf{Fase de Preintegración:} Realización de pruebas CRUD y personalizaciones avanzadas de lógica de negocio por parte del equipo técnico. %
\end{enumerate}

\subsection{Figuras y Evidencias: Estudio Experimental Comparativo}
\markboth{CAPÍTULO 2\quad FIGURAS Y EVIDENCIAS}{}

Para validar la efectividad de \textit{FSMicroGenerator}, Khalfaoui et al. (2025) ejecutaron un estudio experimental controlado comparando dos enfoques de desarrollo para un sistema distribuido de gestión de inventarios bajo un límite estricto de \textbf{8 horas}: %
\begin{itemize}
    \item \textbf{Desarrollo Manual:} 5 desarrolladores experimentados (2-3 años de experiencia en pila MEAN y microservicios). %
    \item \textbf{Desarrollo con FSMicroGenerator:} 2 desarrolladores junior (1 y 2 años de experiencia). %
\end{itemize}

\textoside{
Dado que el desarrollo manual no logró completar el sistema en el tiempo límite, se aplicaron los coeficientes de cobertura funcional ($R_f = 0.61$) y no funcional ($R_{nfr} = 0.33$) para proyectar equitativamente los métricos a alcance completo mediante el factor $R_f \times R_{nfr} = 0.2013$. %
}{CONCEPTO CLAVE\\[0.1cm]Normalización de Cobertura}
\vspace{0.5cm}

\subsubsection{Evaluación Cuantitativa de Eficiencia y Rendimiento}

\begin{table}[htpb!]
    \centering
    \footnotesize
    \renewcommand{\arraystretch}{1.5}
    \begin{tabularx}{\textwidth}{>{\raggedright\arraybackslash\bfseries}p{2.5cm} >{\centering\arraybackslash}p{2.5cm} >{\centering\arraybackslash}p{2.5cm} >{\centering\arraybackslash}p{2.5cm} >{\raggedright\arraybackslash}X}
    \toprule
    Métrica de Evaluación & Desarrollo Manual (Bruto) & Desarrollo Manual (Normalizado) & FSMicroGenerator (Junior) & Impacto / Ventaja \\
    \midrule
    Tiempo de Desarrollo & 8h (Incompleto) \textit{[35]} & $\sim$40h (Proyectado) \textit{[36]} & \textbf{3 horas (100\% completo)} \textit{[35]} & Reducción del 92.5\% de tiempo \textit{[36]} \\
    Cobertura Funcional ($R_f$) & 61\% \textit{[34]} & 100\% \textit{[32]} & \textbf{100\%} \textit{[34]} & Cobertura funcional total \textit{[34]} \\
    Cobertura No Funcional ($R_{nfr}$) & 33\% \textit{[34]} & 100\% \textit{[32]} & \textbf{100\%} \textit{[34]} & Cobertura NFR nativa \textit{[34]} \\
    Commits en Git & 115 commits \textit{[35]} & $\sim$571 commits \textit{[36]} & \textbf{11 commits} \textit{[35]} & Estructura estable sin \textit{code churn} \textit{[35, 38]} \\
    Intervenciones Manuales & 1,305 mods. \textit{[35]} & >25,000 mods. \textit{[36]} & \textbf{148 mods.} \textit{[35]} & Automatización masiva \textit{[36]} \\
    Rendimiento JMeter & 25.44 req/s \textit{[37]} & 5.12 req/s \textit{[37]} & \textbf{21.75 req/s} \textit{[37]} & +324\% capacidad proc. \textit{[37]} \\
    Tasa de Error & 0.0\% \textit{[37]} & 0.0\% \textit{[37]} & \textbf{0.0\%} \textit{[37, 39]} & Estabilidad operativa \textit{[37]} \\
    \bottomrule
    \end{tabularx}
\end{table}

\subsubsection{Evaluación Cualitativa de Arquitectura (Escala Likert 1-5, ISO/IEC 25010)}

Los evaluadores independientes asignaron la máxima calificación (\textbf{5/5}) a la solución generada por \textit{FSMicroGenerator} frente a un promedio de (\textbf{2/5} y \textbf{3/5}) en el desarrollo manual en todas las dimensiones: %
\begin{enumerate}
    \item \textbf{Modularidad Arquitectónica (5/5):} Separación perfecta entre \textit{frontend}, \textit{gateway} y microservicios independientes alineados a dominios del negocio. %
    \item \textbf{Reusabilidad de Código (5/5):} Reutilización extensiva mediante componentes compartidos, interceptores y directivas genéricas. %
    \item \textbf{Claridad y Estándares (5/5):} Consistencia total en la convención de nombres, estructura interna MVC y especificación OpenAPI/Swagger. %
    \item \textbf{Mantenibilidad (5/5):} Pruebas unitarias/integración integradas por cada microservicio, reduciendo dramáticamente el riesgo de regresiones. %
\end{enumerate}

\subsection{Caso Práctico Aplicado: Despliegue de un Microservicio de Gestión Científico}

\subsubsection*{Contexto}
Despliegue del ecosistema \textbf{Smart DLSP} (\textit{Smart Digital Logistic Services Providers}), un proyecto de innovación nacional en Marruecos ejecutado por Smart Automation Technologies para modernizar la gestión logística y de carga. %[cite: 5]

\subsubsection*{Análisis}
Se utilizó \textit{FSMicroGenerator} para modelar la arquitectura completa mediante diagramas de clases UML, generando automáticamente microservicios en Node.js/Express.js, interfaz en Angular y orquestación en Docker Compose. %[cite: 5, 15] El despliegue se estructuró operacionalizando rápidamente el sistema base, e integrando progresivamente agentes inteligentes de IA. %[cite: 5]

\subsubsection*{Interpretación}
El uso de \textit{FSMicroGenerator} garantizó la consistencia arquitectónica entre servicios heterogéneos y redujo el tiempo de desarrollo de semanas a días. %[cite: 5] La arquitectura generada demostró una alta modularidad que permitió incorporar agentes de inteligencia artificial sin alterar los componentes existentes ni generar deuda técnica. %[cite: 5]

% --- NOTAS CAP 2 ---
\newpage
\markboth{CAPÍTULO 2\quad NOTAS BIBLIOGRÁFICAS E HISTÓRICAS}{}
\vspace*{0.8cm}
\noindent{\Huge \textbf{Notas bibliográficas\\[0.2cm] e históricas}}\\[1cm]

La literatura sobre generación automática de código incluye herramientas pioneras como \textit{xGenerator} (Paolone et al., 2020), enfocada en aplicaciones web MVC en Java/Hibernate de carácter monolítico, y plataformas \textit{Backend as a Service} (BaaS) como la de Prasanthan et al. (2023), centradas en la generación de servicios \textit{backend} en Node.js a partir de esquemas JSON. % Asimismo, Tesoriero et al. (2022) propusieron arquitecturas basadas en MDA para pilas PHP/Polymer. %

\textit{FSMicroGenerator} (Khalfaoui et al., 2025) supera estas aproximaciones al integrar de manera nativa la arquitectura distribuida de microservicios, la separación estricta \textit{frontend/backend} bajo el patrón MEAN stack, la observabilidad con Prometheus/Grafana y la automatización DevSecOps alineada a los estándares OMG y métricas de calidad ISO/IEC 25010. %[cite: 1, 16]

% --- RESUMEN CAP 2 ---
\newpage
\markboth{CAPÍTULO 2\quad RESUMEN}{}
\vspace*{0.8cm}
{\Huge \textbf{Resumen del capítulo 2}}
\vspace{1cm}

\begin{resumenlist}
    \item \textbf{Superación del Monolito:} Las arquitecturas monolíticas limitan la escalabilidad, mientras que \textit{FSMicroGenerator} automatiza la transición a microservicios. %[cite: 1, 2]
    \item \textbf{Jerarquía DDD y Plantillas:} Organiza el código en Paquetes, Módulos y Proyectos combinando plantillas estáticas y dinámicas. %[cite: 8, 10, 11]
    \item \textbf{Eficiencia Comprobada:} Logró un 100\% de cobertura en solo 3 horas frente a un desarrollo manual incompleto (61\%) en 8 horas. %
    \item \textbf{Calidad y Seguridad:} Obtuvo la calificación máxima (5/5) en modularidad, claridad y mantenibilidad, incorporando prácticas DevSecOps nativas. %
\end{resumenlist}

\vspace{1cm}
\begin{cajabeige}
\textbf{Síntesis:} La generación automatizada basada en UML elimina la brecha entre conceptualización y despliegue escalable.
\end{cajabeige}

% --- CONCLUSIONES CAP 2 ---
\newpage
\markboth{CAPÍTULO 2\quad CONCLUSIONES}{}
\vspace*{0.8cm}
{\Huge \textbf{Conclusiones del capítulo 2}}
\vspace{1cm}

\subsubsection*{Conclusión principal}
El marco \textit{FSMicroGenerator} demuestra que es posible automatizar la creación de soluciones \textit{full-stack} basadas en microservicios a partir de modelos UML, permitiendo a equipos de investigación desplegar arquitecturas altamente escalables, seguras y mantenibles en una fracción del tiempo tradicional. %[cite: 1]

\subsubsection*{Aportes y límites}
\textbf{Aportes:} Reducción drástica del tiempo de desarrollo (-92.5\%), eliminación del código repetitivo y adopción nativa de DevSecOps. %[cite: 16] \textbf{Límites:} La fase de definición inicial de entidades se realiza de forma semimanual a través de la interfaz visual, requiriendo tiempo de configuración en grandes proyectos. %

\subsubsection*{Preguntas abiertas}
¿Cómo pueden integrarse los modelos probabilísticos de IA con generadores deterministas para refinar diagramas UML mediante diálogo conversacional?
% ==========================================
% CAPÍTULO 3
% ==========================================
\newpage
\section[Vibe Coding y Gobernanza del Código Sintético]{Vibe Coding, Inteligencia Artificial Generativa y Gobernanza del Código Sintético}
\markboth{CAPÍTULO 3\quad VIBE CODING Y GOBERNANZA}{}

\begin{cajabeige}
\textbf{Pregunta Central del Capítulo:} ¿Cuáles son las oportunidades, riesgos técnicos y requisitos de gobernanza asociados a la adopción del Vibe Coding e Inteligencia Artificial Generativa en la creación de software, y cómo se puede prevenir la acumulación de deuda técnica y vulnerabilidades de seguridad en el prototipado científico?\\[0.3cm]
\textbf{Objetivo de Aprendizaje:} El lector comprenderá el concepto de Vibe Coding basándose en revisiones sistemáticas de literatura (PRISMA), comparará su rendimiento frente a métodos manuales y plataformas Low-Code, identificará las patologías del código sintético (Cuello de Botella de Verificación y Arquitecturas de Error) y aplicará el marco de gobernanza VIBE4M junto con análisis estático (SAST).
\end{cajabeige}

\textoside{
El \textit{Vibe Coding} es un paradigma emergente de creación de software en el cual el desarrollador o investigador no escribe código sintáctico determinista directamente, sino que interactúa en lenguaje natural mediante diálogo conversacional con asistentes de Inteligencia Artificial Generativa basados en modelos de lenguaje de gran escala (LLMs).\\[0.5cm]
La mediación de intención desplaza la actividad principal del programador desde la implementación manual hacia la especificación clara de requerimientos, la iteración de prompts y la verificación crítica de los resultados producidos.
}{CONCEPTO CLAVE\\[0.1cm]Mediación de Intención y Vibe Coding}
\vspace{0.5cm}

\subsection{Emergencia del Vibe Coding: Revisión Sistemática de Literatura (SLR PRISMA sobre 61 Estudios)}

El avance vertiginoso de los modelos de lenguaje de gran escala (LLMs) especializados en código —tales como OpenAI Codex, Claude 3.5 Sonnet, GitHub Copilot y plataformas dedicadas como Cursor AI y Claude Code— ha dado origen a una transformación radical conocida popularmente como Vibe Coding.

Para formalizar académicamente este fenómeno, Jauhari, Fathullah, Permana y Nugraha (2026) ejecutaron una Revisión Sistemática de la Literatura (Systematic Literature Review, SLR) aplicando rigurosamente los estándares del protocolo PRISMA (Preferred Reporting Items for Systematic Reviews and Meta-Analyses). Tras una búsqueda comprehensiva en bases de datos científicas internacionales (IEEE Xplore, ACM Digital Library, ScienceDirect, Scopus y arXiv) y un filtrado metodológico estricto, seleccionaron y analizaron 61 estudios académicos arbitrados publicados entre 2023 y 2026.

Los hallazgos del estudio de Jauhari et al. (2026) confirman que el Vibe Coding representa la interfaz de desarrollo de menor fricción sintáctica en la historia de la informática. El investigador expresa la vibra o intención funcional de la aplicación en lenguaje natural (por ejemplo: "Crea una interfaz en React que capture archivos CSV de datos espectroscópicos, limpie los valores nulos y los grafique en una curva interactiva con D3.js"), y el agente de IA genera instantáneamente el código fuente, la configuración de estilos y los componentes requeridos.

\subsection{Análisis Comparativo de Rendimiento: Desarrollo Manual vs. Low-Code vs. Vibe Coding con LLMs}

A partir de la síntesis cuantitativa y cualitativa de las 61 investigaciones analizadas por Jauhari et al. (2026), se establece un marco comparativo entre los tres paradigmas dominantes de creación de software:

\begin{enumerate}[label=\textbf{\Alph*.}, leftmargin=1.5cm]
    \item \textbf{Velocidad y Productividad de Prototipado:} La investigación sistemática demuestra de forma consistente que la asistencia por IA en la generación de código permite reducir el tiempo de desarrollo de un Mínimo Viable de Prueba (MVP) entre un 40\% y un 60\% en comparación con el desarrollo manual convencional.
    
    \item \textbf{Nivel de Predictibilidad y Control Arquitectónico:} A diferencia de las plataformas Low-Code (que imponen límites rígidos a través de componentes visuales deterministas que garantizan la consistencia del sistema), el Vibe Coding opera sobre modelos estocásticos probabilísticos. La misma instrucción en lenguaje natural puede producir variaciones en la arquitectura del código generado, exigiendo un mayor nivel de inspección por parte del usuario.
\end{enumerate}

\vspace{0.5cm}
\textbf{C. Matriz Comparativa Multidimensional de Paradigmas}

\begin{table}[htpb!]
    \centering
    \footnotesize
    \renewcommand{\arraystretch}{1.5}
    \begin{tabularx}{\textwidth}{>{\raggedright\arraybackslash\bfseries}p{2.8cm} >{\raggedright\arraybackslash}X >{\raggedright\arraybackslash}X >{\raggedright\arraybackslash}X}
    \toprule
    Dimensión Técnica & \textbf{Desarrollo Manual Tradicional} & \textbf{Low-Code / No-Code Determinista} & \textbf{Vibe Coding / IA Generativa Pura} \\
    \midrule
    Velocidad de Prototipado & Baja a Moderada (Semanas/Meses) & Alta (MVP en Horas/Días) & Extremadamente Alta (Aceleración entre 40\% y 60\%) \\
    Control Arquitectónico & Total y Absoluto (Línea a línea) & Moderado-Alto (Marco visual) & Variable / Probabilístico (Prompt) \\
    Habilidades Requeridas & Sintaxis compleja, algoritmos. & Lógica de negocio, modelado UML. & Ingeniería de prompts, verificación. \\
    Riesgo de Vulnerabilidades & Controlable manualmente. & Bajo (RBAC y gobernanza nativa). & Alto sin SAST (Fallos OWASP). \\
    Flexibilidad de Expresión & Limitada por lenguaje. & Limitada por componentes visuales. & Máxima (Lenguaje natural libre). \\
    Acumulación de Deuda Técnica & Manejable vía refactorización. & Controlada por la plataforma. & Muy Alta si no se aplica VIBE4M. \\
    \bottomrule
    \end{tabularx}
\end{table}

\subsection{Patologías del Código Generado por IA: Cuello de Botella de Verificación, Arquitecturas de Error y Deuda Técnica}

A pesar de su velocidad, la adopción del Vibe Coding sin controles metodológicos introduce patologías técnicas graves:

\begin{enumerate}
    \item \textbf{El Cuello de Botella de Verificación (Verification Bottleneck):} Como destacan Jauhari et al. (2026), la IA no elimina el esfuerzo de desarrollo, sino que desplaza masivamente la carga cognitiva hacia la fase posterior de pruebas e inspección. Se genera un efecto de congestión donde el desarrollador pasa la mayor parte del tiempo intentando leer, comprender, probar y depurar grandes volúmenes de código sintético que no escribió manualmente.
    \item \textbf{Arquitecturas de Error y Crisis de Deuda Técnica:} Investigaciones cuantitativas recientes (van Hurne, 2025; Zhao et al., 2025) demuestran que el código producido tiende a duplicar funciones, omitir el manejo de excepciones de borde (edge cases), ignorar principios SOLID y acumular una elevada deuda técnica.
    \item \textbf{Vulnerabilidades Ocultas de Seguridad:} El estudio de Jauhari et al. (2026) señala que entre el 30\% y el 45\% del código generado sin supervisión por LLMs modernos contiene vulnerabilidades de seguridad clasificadas en el Top 10 de OWASP:
    \begin{itemize}
        \item Inyección SQL y desinfectación deficiente de entradas de usuario.
        \item Uso de claves secretas (hardcoded credentials) o tokens incrustados en texto plano.
        \item Librerías desactualizadas con vulnerabilidades conocidas (CVEs).
        \item Configuración insegura de políticas CORS (Cross-Origin Resource Sharing).
    \end{itemize}
\end{enumerate}

\subsection{Caso Práctico Aplicado: Auditoría de Seguridad y Refactorización}

\subsubsection*{Contexto}
Un prototipo de Flask generado por IA incluyó vulnerabilidades de Inyección SQL directa:

\begin{tcolorbox}[colback=bg_marfil, colframe=unmsm_gold, arc=3pt, boxrule=0.6pt, top=0.2cm, bottom=0.2cm, left=0.4cm, right=0.4cm]
\footnotesize
\begin{verbatim}
# Código generado por IA (Inseguro - Inyección SQL)
@app.route('/get_data')
def get_data():
    filename = request.args.get('filename')
    query = f"SELECT * FROM measurements WHERE file = '{filename}'"
    result = db.engine.execute(query)
    return jsonify([dict(row) for row in result])
\end{verbatim}
\end{tcolorbox}

\subsubsection*{Análisis y Refactorización}
El marco VIBE4M interceptó el fallo vía Bandit (SAST) en CI/CD. El investigador refactorizó el código usando \textit{safe pathing} y consultas parametrizadas:

\begin{tcolorbox}[colback=bg_marfil, colframe=unmsm_gold, arc=3pt, boxrule=0.6pt, top=0.2cm, bottom=0.2cm, left=0.4cm, right=0.4cm]
\footnotesize
\begin{verbatim}
# Código refactorizado bajo gobernanza VIBE4M (Seguro)
@app.route('/get_data')
def get_data():
    filename = request.args.get('filename')
    if not filename or not re.match(r'^[a-zA-Z0-9_\-\.]+\.csv$', filename):
        return jsonify({'error': 'Nombre de archivo inválido'}), 400
    stmt = text("SELECT * FROM measurements WHERE file = :fn")
    result = db.engine.execute(stmt, {'fn': filename})
    return jsonify([dict(row) for row in result])
\end{verbatim}
\end{tcolorbox}

% --- NOTAS CAP 3 ---
\newpage
\markboth{CAPÍTULO 3\quad NOTAS BIBLIOGRÁFICAS E HISTÓRICAS}{}
\vspace*{0.8cm}
\noindent{\Huge \textbf{Notas bibliográficas\\[0.2cm] e históricas}}\\[0.5cm]

El estudio del impacto de los modelos de lenguaje de gran escala (LLMs) en la programación se apoya en los trabajos sobre la arquitectura Transformer de Vaswani et al. (2017) y la alineación por retroalimentación humana (RLHF) de Ouyang et al. (2022). La definición del Vibe Coding y la mediación de intenciones fue formalizada por Meske et al. (2025) en IEEE Access. Las investigaciones sistemáticas sobre oportunidades y riesgos se sostienen en la revisión PRISMA de Jauhari, Fathullah, Permana y Nugraha (2026). Los marcos de prevención de "Arquitecturas de Error" y análisis estático SAST se fundamentan en Sartori (2025), Maes (2025), van Hurne (2025) y Zhao et al. (2025).

% --- RESUMEN CAP 3 ---
\newpage
\markboth{CAPÍTULO 3\quad RESUMEN}{}
\vspace*{0.8cm}
{\Huge \textbf{Resumen del capítulo 3}}
\vspace{1cm}
\begin{resumenlist}
    \item \textbf{Vibe Coding:} Acelera prototipos MVP entre 40-60\%.
    \item \textbf{Patologías:} Produce cuellos de botella de verificación y "arquitecturas de error".
    \item \textbf{Seguridad:} Fallos OWASP comunes requieren supervisión humana obligatoria.
    \item \textbf{VIBE4M:} Marco de gobernanza exige SAST, segmentación y revisión.
\end{resumenlist}

\vspace{1cm}
\begin{cajabeige}
\textbf{Síntesis:} El Vibe Coding requiere marcos de gobernanza rigurosos (VIBE4M y SAST) para evitar que la aceleración sintáctica genere sistemas frágiles, inseguros o difíciles de mantener.
\end{cajabeige}

% --- CONCLUSIONES CAP 3 ---
\newpage
\markboth{CAPÍTULO 3\quad CONCLUSIONES}{}
\vspace*{0.8cm}
{\Huge \textbf{Conclusiones del capítulo 3}}
\vspace{1cm}
\subsubsection*{Conclusión principal}
El Vibe Coding representa la interfaz de desarrollo de software de mayor velocidad en la historia de la informática, permitiendo a investigadores de cualquier disciplina materializar prototipos funcionales mediante lenguaje natural. Sin embargo, su adopción en entornos científicos debe estar acompañada por marcos de gobernanza (VIBE4M) y herramientas de análisis estático (SAST), asegurando que la agilidad no comprometa la calidad técnica ni la seguridad.

\subsubsection*{Aportes y límites}
\textbf{Aportes:} Descomposición del impacto de la IA en las 6 etapas del ciclo SDLC, sistematización del marco VIBE4M y demostración práctica de corrección de vulnerabilidades OWASP. \textbf{Límites:} El análisis SAST no detecta fallas lógicas complejas de dominio; el éxito del Vibe Coding depende de la capacidad del usuario para formular prompts precisos y verificar resultados.

\subsubsection*{Preguntas abiertas}
¿Cómo pueden evolucionar las herramientas de prueba automatizada para realizar verificaciones de corrección semántica en tiempo real mientras el investigador interactúa por voz o chat con un agente generador de código?

% =======% ==========================================
% CAPÍTULO 4
% ==========================================
\newpage
\section[La Escalera de Capacidades y Supervisión con IA]{La Escalera de Capacidades y la Supervisión de Sistemas con Inteligencia Artificial}
\markboth{CAPÍTULO 4\quad ESCALERA DE CAPACIDADES}{} %[cite: 18]

\begin{cajabeige}
\textbf{Pregunta Central del Capítulo:} ¿Cómo permite la Escalera de Capacidades describir la autonomía de los sistemas y organizar las evidencias necesarias para supervisarlos responsablemente? \\[0.3cm] %[cite: 18]
\textbf{Objetivo de Aprendizaje:} Distinguir los niveles L0–L4, identificar los componentes de la capa de aseguramiento y analizar las posibilidades y limitaciones de su aplicación en educación e investigación. %[cite: 18]
\end{cajabeige}

\textoside{
La incorporación de inteligencia artificial puede cambiar la distribución de las tareas dentro de una actividad profesional. %[cite: 18] Algunas operaciones rutinarias pueden automatizarse, mientras que adquieren mayor importancia la formulación del problema, la verificación de resultados, el diseño del sistema y la supervisión de sus acciones. %[cite: 18] Memari y Rudolph (2026) utilizan esta idea para proponer una actualización de la formación informática y profesional. %[cite: 18] Su planteamiento no implica que todos los empleos estén protegidos frente a la automatización ni que sus efectos sean iguales en todas las actividades. %[cite: 18]
}{CONCEPTO CLAVE\\[0.1cm]Reasignación de Tareas}
\vspace{0.5cm}

\subsection{Reasignación de Tareas y Transformación del Trabajo} %[cite: 18]

Las herramientas de IA pueden acelerar la elaboración de código, documentos y primeras versiones de una solución. %[cite: 18] Sin embargo, obtener rápidamente un resultado no demuestra que sea correcto, seguro o adecuado para el problema. %[cite: 18]

Por ello, la formación debe desarrollar capacidades para definir requisitos, detectar errores, justificar decisiones y comprobar el comportamiento de los sistemas. %[cite: 18] La velocidad de producción debe acompañarse de evidencias de comprensión y verificación. %[cite: 18]

Memari y Rudolph (2026, secciones 2–4) fundamentan su propuesta en una revisión narrativa de literatura sobre trabajo, ingeniería de software y educación. %[cite: 18] Los autores reconocen que los cambios laborales también están relacionados con factores económicos y organizativos; por tanto, no atribuyen todos esos cambios exclusivamente a la IA. %[cite: 18]

Para una guía de prototipado científico, este enfoque permite formular una distinción metodológica: construir software que funciona y validar una hipótesis científica son actividades relacionadas, pero requieren evidencias diferentes. %[cite: 18] El funcionamiento del prototipo debe comprobarse frente a sus requisitos; las conclusiones científicas necesitan, además, un método de investigación adecuado. %[cite: 18] La Escalera de Capacidades no sustituye esa evaluación. %[cite: 18]

\subsection{Los Cinco Niveles de la Escalera de Capacidades} %[cite: 18]

La Escalera de Capacidades clasifica los sistemas según su autonomía operativa y las responsabilidades humanas necesarias para supervisarlos (Memari y Rudolph, 2026, sección 5). %[cite: 18]

\begin{table}[htpb!]
    \centering
    \footnotesize
    \renewcommand{\arraystretch}{1.5}
    \begin{tabularx}{\textwidth}{>{\raggedright\arraybackslash\bfseries}p{2.5cm} >{\raggedright\arraybackslash}X >{\raggedright\arraybackslash}X}
    \toprule
    Nivel & \textbf{Característica principal} & \textbf{Evidencias para evaluar su funcionamiento} \\ %[cite: 18]
    \midrule
    L0. Disparador & Un evento inicia un proceso: envío de un formulario, programación horaria o recepción de un webhook. & Descripción del evento, entradas y explicación de cuándo se activa. \\ %[cite: 18]
    L1. Automatización & Una acción se ejecuta a partir de un disparador. & Regla de automatización, caso de prueba y resultado observado. \\ %[cite: 18]
    L2. Flujo de trabajo & Un proceso fijo encadena varios pasos entre herramientas: recopilar, validar, enviar y registrar información. & Diagrama del flujo, reglas de validación, pruebas y manejo de errores. \\ %[cite: 18]
    L3. Agente de IA & Un sistema orientado a una meta planifica y actúa bajo verificación humana. & Meta, configuración, muestras de resultados, lista de verificación y punto de aprobación humana. \\ %[cite: 18]
    L4. Equipo de agentes & Varios agentes distribuyen tareas y presentan resultados bajo supervisión humana. & Mapa de funciones, estructura de delegación, registro de auditoría, análisis de riesgos e informe de decisión humana. \\ %[cite: 18]
    \bottomrule
    \end{tabularx}
    \caption*{Fuente de la tabla: adaptación de Memari y Rudolph (2026, figura 1 y tabla 2).} %[cite: 18]
\end{table}

Los niveles superiores requieren mayor responsabilidad de supervisión. %[cite: 18] No representan una escala de inteligencia, calidad científica o superioridad de una solución. %[cite: 18] Tampoco constituyen una secuencia que todos los estudiantes o proyectos deban completar. %[cite: 18]

La selección del nivel debe responder al problema y a la capacidad de verificar el sistema. %[cite: 18] Incorporar más autonomía no constituye, por sí mismo, una mejora demostrada. %[cite: 18]

\subsection{Capa Transversal de Aseguramiento y Evaluación} %[cite: 18]

El marco incorpora cinco componentes de aseguramiento aplicables en todos los niveles (Memari y Rudolph, 2026, sección 5): %[cite: 18]

\begin{enumerate}
    \item \textbf{Verificación:} comprobar los resultados mediante pruebas, revisión y detección de errores. %[cite: 18]
    \item \textbf{Seguridad:} identificar riesgos, proteger información y controlar las acciones autorizadas. %[cite: 18]
    \item \textbf{Documentación:} registrar configuraciones, decisiones, pruebas y uso de herramientas. %[cite: 18]
    \item \textbf{Ética:} considerar atribución, privacidad, posibles daños y cumplimiento de las políticas aplicables. %[cite: 18]
    \item \textbf{Aprobación humana:} establecer quién revisa y autoriza los resultados o acciones que requieren supervisión. %[cite: 18]
\end{enumerate}

Estos componentes deben concretarse según las características y los riesgos del sistema. %[cite: 18] El artículo no establece una obligación universal de realizar análisis estático de código, calcular porcentajes de edición o firmar formalmente todas las aprobaciones. %[cite: 18] Esas medidas pueden incorporarse cuando el contexto las justifique. %[cite: 18]

\subsubsection*{Divulgación del uso de IA} %[cite: 18]

Memari y Rudolph (2026, sección 7) proponen una plantilla breve que permita responder: %[cite: 18]

\begin{itemize}
    \item ¿Qué herramienta de IA se utilizó? %[cite: 18]
    \item ¿Qué se le pidió? %[cite: 18]
    \item ¿Qué partes de la respuesta se aceptaron, rechazaron o modificaron? %[cite: 18]
    \item ¿Cómo se verificó el resultado? %[cite: 18]
    \item ¿Qué comprende ahora el estudiante que antes no comprendía? %[cite: 18]
\end{itemize}

La última pregunta permite relacionar el uso de la herramienta con el aprendizaje. %[cite: 18] Un registro de instrucciones, por sí solo, no demuestra comprensión. %[cite: 18]

\subsubsection*{Evaluación de la comprensión} %[cite: 18]

La fuente combina pruebas, entregas parciales, portafolios, registros y defensas orales. %[cite: 18] Propone ajustar el esfuerzo de evaluación al riesgo y a la carga docente; como ejemplo, incluye defensas breves de dos a tres minutos para una muestra del 20–40 \% del alumnado (Memari y Rudolph, 2026, sección 7 y tabla 4). %[cite: 18]

Estas estrategias aportan evidencias sobre el proceso y la autoría, pero no garantizan por sí solas la ausencia de plagio ni demuestran que se haya evitado la pérdida de habilidades. %[cite: 18]

\subsection{Experiencia Educativa Documentada y Evidencia Experimental Complementaria} %[cite: 18]

\subsubsection*{Experiencia piloto descrita por Memari y Rudolph} %[cite: 18]

Memari y Rudolph (2026, sección 9 y tabla 5) describen una asignatura optativa interdisciplinaria desarrollada durante dos semestres: otoño de 2025 y primavera de 2026. %[cite: 18] Participaron aproximadamente 40 estudiantes por edición, unos 80 en total, procedentes de computación y negocios. %[cite: 18] La asignatura no exigía conocimientos previos de programación. %[cite: 18]

Los estudiantes formaron equipos con funciones diferenciadas, utilizaron una plataforma comercial de orquestación de agentes sin código y abordaron problemas aportados por socios industriales regionales. %[cite: 18] La evaluación se centró en evidencias correspondientes a cada función: portafolios, demostraciones, reflexiones y verificación de resultados. %[cite: 18]

Los autores observaron que estudiantes de negocios podían trabajar con agentes y equipos de agentes sin haber cursado previamente programación. %[cite: 18] El artículo sitúa la experiencia en una universidad pública orientada a la docencia. %[cite: 18]

\subsubsection*{Alcance de la evidencia del piloto} %[cite: 18]

La experiencia es exploratoria, corresponde a un solo programa y no dispone de grupo de control. %[cite: 18] El seguimiento posterior fue informal y se realizó con una minoría autoseleccionada de participantes. %[cite: 18] Por ello, permite identificar posibilidades de aplicación y preguntas para futuras investigaciones, pero no demostrar efectos causales sobre aprendizaje o empleabilidad. %[cite: 18]

\subsubsection*{Evidencia experimental complementaria: asistencia guiada en programación} %[cite: 18]

Noraset et al. (2026a) realizaron un experimento controlado con 42 participantes en un curso breve de programación en Python. %[cite: 18] Compararon dos diseños de chatbot: uno que proporcionaba orientación sin entregar soluciones directas y otro que permitía un uso sin restricciones. %[cite: 18]

El grupo que utilizó el asistente guiado presentó una mejora significativamente mayor entre la evaluación inicial y la final. %[cite: 18] Este resultado aporta evidencia sobre la importancia del diseño de la interacción con IA para apoyar el aprendizaje. %[cite: 18] Su alcance corresponde a las condiciones del curso estudiado; no demuestra efectos universales o duraderos. %[cite: 18]

El estudio complementa el capítulo al documentar una experiencia evaluada en la que se buscó preservar la participación del estudiante durante la resolución de problemas. %[cite: 18] No constituye una validación de la Escalera de Capacidades ni una evaluación del prototipado científico. %[cite: 18] Su contribución es respaldar, en un contexto educativo concreto, la discusión sobre cómo organizar la asistencia de IA sin sustituir completamente el trabajo del usuario. %[cite: 18]

La fe de erratas del estudio identifica a Mahidol University como institución donde se realizó la investigación y completa información sobre sus recursos y aprobación ética (Noraset et al., 2026b). %[cite: 18] El caso se atribuye a sus autores y no se presenta como una prueba realizada por el autor de esta guía. %[cite: 18]

% --- NOTAS CAP 4 ---
\newpage
\markboth{CAPÍTULO 4\quad NOTAS BIBLIOGRÁFICAS E HISTÓRICAS}{}
\vspace*{0.8cm}
\noindent{\Huge \textbf{Notas bibliográficas\\[0.2cm] e históricas}}\\[1cm]

La clasificación L0–L4 y la capa de aseguramiento proceden de la sección 5 de Memari y Rudolph (2026). %[cite: 18] Las estrategias de evaluación se desarrollan en la sección 7; el piloto, en la sección 9; y las limitaciones, en la sección 11. %[cite: 18] El documento está disponible como prepublicación en arXiv. %[cite: 18]

El trabajo de Noraset et al. (2026a), publicado en Computers and Education: Artificial Intelligence, aporta evidencia experimental complementaria sobre aprendizaje con asistencia guiada. %[cite: 18] Debe consultarse junto con su fe de erratas (2026b). %[cite: 18] El artículo estuvo disponible en línea en diciembre de 2025 y pertenece al volumen de junio de 2026. %[cite: 18]

Las dos investigaciones cumplen funciones diferentes: una propone un marco e ilustra su aplicación educativa; la otra compara diseños de asistencia en una tarea formativa. %[cite: 18] No deben combinarse sus resultados como si constituyeran una validación experimental única del marco. %[cite: 18]

% --- RESUMEN CAP 4 ---
\newpage
\markboth{CAPÍTULO 4\quad RESUMEN}{}
\vspace*{0.8cm}
{\Huge \textbf{Resumen del capítulo 4}}
\vspace{1cm}

\begin{resumenlist}
    \item La automatización puede redistribuir tareas y aumentar la importancia de la verificación y la supervisión. %[cite: 18]
    \item Los niveles L0–L4 describen autonomía operativa y responsabilidades humanas. %[cite: 18]
    \item El aseguramiento comprende verificación, seguridad, documentación, ética y aprobación humana. %[cite: 18]
    \item El piloto de Memari y Rudolph ilustra una aplicación educativa con evidencia exploratoria. %[cite: 18]
    \item El experimento de Noraset et al. aporta evidencia sobre aprendizaje con asistencia guiada en un curso breve de Python. %[cite: 18]
    \item Ninguno de esos resultados demuestra por sí solo la validez científica de un prototipo o una garantía general de efectividad de la Escalera de Capacidades. %[cite: 18]
\end{resumenlist}

\vspace{1cm}
\begin{cajabeige}
\textbf{Síntesis:} La Escalera de Capacidades aporta una estructura para decidir qué tareas delegar y qué evidencias documentar, pero su aplicación debe evaluarse siempre en función de su contexto y no garantiza automáticamente el rigor científico. %[cite: 18]
\end{cajabeige}

% --- CONCLUSIONES CAP 4 ---
\newpage
\markboth{CAPÍTULO 4\quad CONCLUSIONES}{}
\vspace*{0.8cm}
{\Huge \textbf{Conclusiones del capítulo 4}}
\vspace{1cm}

\subsubsection*{Conclusión principal}
La Escalera de Capacidades ofrece una estructura para decidir qué tareas delegar, qué resultados verificar y qué evidencias documentar. %[cite: 18] Puede orientar el diseño y la evaluación de sistemas con IA, siempre que su aplicación se ajuste al contexto. %[cite: 18]

La experiencia educativa exploratoria y el experimento sobre asistencia guiada aportan antecedentes de diferente alcance. %[cite: 18] Respaldan la discusión sobre supervisión y diseño de la interacción, pero no permiten afirmar que el marco garantice validez científica, integridad académica o protección frente a la dependencia excesiva de la IA. %[cite: 18]

\subsubsection*{Aportes y límites}
El aporte principal del marco es vincular clases de sistemas con responsabilidades y evidencias de supervisión. %[cite: 18] Entre las limitaciones de su fundamentación se encuentran el carácter narrativo de la revisión, el piloto de una sola institución y la ausencia de seguimiento sistemático. %[cite: 18] Se necesita validación en más cursos, instituciones y poblaciones. %[cite: 18]

El experimento complementario se limita a 42 participantes y a un curso breve de Python. %[cite: 18] Sus resultados deben interpretarse dentro de ese contexto; trasladarlos al prototipado científico requiere evidencia adicional. %[cite: 18]

\subsubsection*{Pregunta abierta}
¿Qué combinación de pruebas, documentación y revisión humana permite evaluar un prototipo con suficiente rigor y una carga de trabajo sostenible? %[cite: 18]

\subsubsection*{Referencias del Capítulo 4} %[cite: 18]
\begin{itemize}
    \item Memari, M., \& Rudolph, G. (2026). The Capability Ladder: A Curriculum-Modernization Framework for Workforce Readiness in the AI Era [Prepublicación]. arXiv. \url{https://doi.org/10.48550/arXiv.2608.07779} %[cite: 18]
    \item Noraset, T., Supratak, A., Ragkhitwetsagul, C., Worathong, N., \& Tuarob, S. (2026a). Evaluating lab assistant chatbot on student learning and behaviors in a programming short course. \textit{Computers and Education: Artificial Intelligence}, 10, 100527. \url{https://doi.org/10.1016/j.caeai.2025.100527} %[cite: 18]
    \item Noraset, T., Supratak, A., Ragkhitwetsagul, C., Worathong, N., \& Tuarob, S. (2026b). Corrigendum to “Evaluating lab assistant chatbot on student learning and behaviors in a programming short course”. \textit{Computers and Education: Artificial Intelligence}, 10, 100547. \url{https://doi.org/10.1016/j.caeai.2026.100547} %[cite: 18]
\end{itemize}
\newpage
\section[Síntesis Metodológica y Hoja de Ruta Institucional]{Síntesis Metodológica, Matriz de Decisión y Hoja de Ruta para la Validación Científica}
\markboth{CAPÍTULO 5\quad SÍNTESIS Y HOJA DE RUTA}{}

\begin{cajabeige}
\textbf{Pregunta Central del Capítulo:} ¿Cómo integrar de manera armónica plataformas Low-Code deterministas, generadores basados en modelos UML y Vibe Coding asistido por IA en una hoja de ruta institucional que maximice la velocidad de validación empírica sin comprometer la gobernanza, la seguridad ni la calidad técnica?\\[0.3cm]
\textbf{Objetivo de Aprendizaje:} Sintetizar las fortalezas comparativas de cada paradigma, aplicar matrices de decisión arquitectónica multidimensionales y estructurar un plan maestro de adopción tecnológica y aseguramiento de calidad en laboratorios de investigación y educación superior.
\end{cajabeige}

\textoside{
Consiste en integrar deliberadamente la herramienta computacional óptima según el componente del sistema: bajo código determinista para flujos documentales y gobernanza; generación guiada por UML para arquitecturas de backend y microservicios; y Vibe Coding con IA generativa para maquetación ágil e interfaces de usuario.
}{CONCEPTO CLAVE\\[0.1cm]Orquestación Híbrida}
\vspace{0.5cm}

\subsection{Sostenibilidad Metodológica y Paradigma de Orquestación Híbrida}

La aceleración en la validación empírica de hipótesis científicas no se logra adoptando una única herramienta tecnológica de forma aislada, sino mediante el establecimiento de un modelo de \textbf{Orquestación Híbrida}. Como demuestran las investigaciones de Marian et al. (2026), Khalfaoui et al. (2025) y Jauhari et al. (2026), los tres paradigmas emergentes —Low-Code empresarial determinista, generación basada en modelos UML e Inteligencia Artificial Generativa (\textit{Vibe Coding})— presentan fortalezas complementarias y limitaciones cruzadas.

El principio fundamental de la Orquestación Híbrida radica en desacoplar el prototipo científico en capas funcionales y asignar a cada capa el motor de construcción que maximice la eficiencia y minimice los riesgos de producción:

\begin{enumerate}
    \item \textbf{Capa de Gobernanza y Procesos Documentales:} Asignada a plataformas Low-Code deterministas (ej. Microsoft Power Platform), aprovechando la autenticación centralizada mediante Entra ID, la persistencia en repositorios estructurados y el cumplimiento nativo de normas de auditoría documental (ISO 15489 e ISO 27001).
    \item \textbf{Capa de Backend y Arquitectura de Datos Escalable:} Asignada a generadores dirigidos por modelos (ej. FSMicroGenerator), transformando diagramas de clases UML en arquitecturas de microservicios \textit{full-stack} (stack MEAN) empaquetadas en contenedores Docker y monitoreadas en tiempo real.
    \item \textbf{Capa de Exploración Inicial de UI/UX e Interacción en Lenguaje Natural:} Asignada al \textit{Vibe Coding} respaldado por modelos de lenguaje (LLMs), permitiendo a los investigadores maquetar interfaces y probar conceptos de usuario en ciclos de horas, bajo la supervisión de herramientas de análisis estático de código (SAST).
\end{enumerate}

Este enfoque consolida el paradigma de \textbf{reasignación de tareas} (\textit{task reallocation}) teorizado por Memari y Rudolph (2026), garantizando que la automatización asuma la sobrecarga sintáctica mientras los investigadores ejercen el control arquitectónico y la validación de resultados.

\subsection{Matriz Multidimensional Unificada de Selección Tecnológica}

Para guiar la toma de decisiones técnicas en proyectos de I+D, se presenta la evaluación comparativa multidimensional consolidada a partir de los datos empíricos de las investigaciones analizadas:

\vspace{0.5cm}
\begin{center}
    \footnotesize
    \renewcommand{\arraystretch}{1.5}
    \begin{tabularx}{\textwidth}{>{\raggedright\arraybackslash\bfseries}p{2.8cm} >{\raggedright\arraybackslash}X >{\raggedright\arraybackslash}X >{\raggedright\arraybackslash}X}
    \toprule
    Dimensión de Análisis & \textbf{Low-Code Empresarial (Power Platform)} & \textbf{Generación UML (FSMicroGenerator)} & \textbf{Vibe Coding asistido por IA (LLMs)} \\
    \midrule
    Propósito Arquitectónico & Automatización de flujos documentales y aprobaciones. & Backend distribuido y microservicios escalables. & Maquetación rápida de UI/UX y exploración de conceptos. \\
    Velocidad de Prototipado & Alta (Reducción de tiempo de ciclo del 65.6\%). & Extremadamente Alta (100\% funcionalidad en 3 horas). & Ultra Rápida (Aceleración de MVP entre 40\% y 60\%). \\
    Control Arquitectónico & Moderado-Alto (Restringido por políticas del proveedor). & Total (Basado en patrones DDD y modelos UML). & Bajo-Moderado (Dependiente de la precisión del prompt). \\
    Gobernanza y Auditoría & Nativa (Auditoría inmutable e ISO 15489). & Monitoreo expuesto en Prometheus y Grafana. & Requiere implementación externa de SAST / VIBE4M. \\
    Esfuerzo Infraestructural & Mínimo (Nube SaaS/PaaS gestionada). & Moderado (Despliegue de contenedores Docker/GitLab). & Variable (Depende del entorno de ejecución del código). \\
    Perfil de Usuario & Desarrollador Ciudadano / Investigador de campo. & Ingeniero de Software / Arquitecto de Datos. & Investigador multidisciplinario / Prototipador. \\
    \bottomrule
    \end{tabularx}
\end{center}
\vspace{0.5cm}

\subsection{Análisis de Stakeholders (Matriz Poder-Interés) y Evaluación de TCO}

La implementación de prototipos en organizaciones académicas y centros de investigación requiere gestionar las expectativas de los actores clave mediante una \textbf{Matriz de Poder e Interés}:

\begin{itemize}
    \item \textbf{Gestionar Atentamente (Alto Poder, Alto Interés):} Investigadores Principales (PI) y Directores de Laboratorio. Exigen validez científica, rigor en el procesamiento de datos y cumplimiento de los plazos del proyecto.
    \item \textbf{Mantener Satisfechos (Alto Poder, Bajo Interés):} Comités de Ética, Oficiales de Seguridad de la Información (CISO) y Auditores Institucionales. Su prioridad es la protección de datos sensibles, el cumplimiento de la norma ISO 27001 y el control de accesos basado en roles (RBAC).
    \item \textbf{Mantener Informados (Bajo Poder, Alto Interés):} Estudiantes de posgrado, asistentes de investigación y desarrolladores ciudadanos. Requieren capacitación en herramientas de bajo código y guías transparentes sobre el uso ético de IA.
    \item \textbf{Monitorear (Bajo Poder, Bajo Interés):} Usuarios finales externos y evaluadores de pruebas piloto.
\end{itemize}

\subsection{Evaluación del Costo Total de Propiedad (TCO)}

El análisis económico revela diferencias estructurales entre los enfoques adoptados para el prototipado:

\begin{itemize}
    \item \textbf{Enfoque Low-Code SaaS/PaaS (Power Platform):} Optimiza el TCO al aprovechar licencias educativas o corporativas preexistentes (Microsoft 365). La utilización de listas de SharePoint con columnas de búsqueda (\textit{Lookup Columns}) permite modelar relaciones de datos sin incurrir en costos adicionales de bases de datos dedicadas como Microsoft Dataverse.
    \item \textbf{Enfoque Open-Source basado en Modelos (FSMicroGenerator):} Minimiza los costos recurrentes de licenciamiento al generar código ejecutable sobre tecnologías de código abierto (MongoDB, Express, Angular, Node.js, RabbitMQ). Sin embargo, requiere inversión inicial en infraestructura de cómputo para hospedar contenedores Docker y servidores de integración continua.
\end{itemize}

\subsection{Algoritmo Formulado y Reglas de Decisión Arquitectónica}

Para sistematizar la elección de herramientas en el laboratorio, se formula el siguiente algoritmo de decisión lógica:

\begin{tcolorbox}[colback=bg_marfil, colframe=unmsm_gold, arc=3pt, boxrule=0.6pt, top=0.25cm, bottom=0.25cm, left=0.4cm, right=0.4cm]
\footnotesize
\begin{verbatim}
INICIO: Evaluación del Proyecto de Investigación
|
+-- ¿El proyecto requiere gestión de documentos, firmas o auditoría ISO?
|   |-- SÍ: Seleccionar Low-Code Empresarial (Power Apps + Automate)
|
+-- ¿El proyecto requiere backend de alta concurrencia o DDD?
|   |-- SÍ: Diseñar UML -> Ejecutar FSMicroGenerator (Stack MEAN)
|
+-- ¿El proyecto busca validar una interfaz de usuario o concepto exploratorio?
    |-- SÍ: Aplicar Vibe Coding con IA Generativa
    |
    +--> APLICAR CAPA DE ASEGURAMIENTO (Obligatorio):
         1. Registrar Prompts de generación.
         2. Ejecutar escaneo SAST (Trivy / SonarQube).
         3. Validar punto de Aprobación Humana (Human-in-the-Loop).
FIN
\end{verbatim}
\end{tcolorbox}

\subsection{Mitigación del Triángulo de Riesgos y Capa de Aseguramiento}

El prototipado acelerado expone a las instituciones a tres riesgos interconectados conocidos como el \textbf{Triángulo de Riesgos}:

\begin{enumerate}
    \item \textbf{Acumulación de Deuda Técnica:} Código generado por IA sintética sin refactorización ni pruebas de integración.
    \item \textbf{Atrofia Cognitiva (\textit{Hollow Senior}):} Dependencia ciega del investigador en las salidas del modelo de lenguaje, perdiendo la capacidad de depurar errores complejos.
    \item \textbf{Vulnerabilidades de Seguridad Ocultas:} Inyección de código propenso a fallos de ejecución o brechas de autenticación.
\end{enumerate}

Para neutralizar estos riesgos, se exige la implementación transversal de la \textbf{Capa de Aseguramiento} (\textit{Assurance Layer}) propuesta por Memari y Rudolph (2026) y reforzada por el marco de gobernanza \textbf{VIBE4M} (Jauhari et al., 2026; Maes, 2025). Esta capa impone cinco mecanismos de auditoría obligatorios:

\begin{enumerate}
    \item \textbf{Plantilla de Divulgación de IA:} Registro detallado donde el investigador documenta los modelos utilizados, los prompts formulados, el código generado aceptado/rechazado y las modificaciones manuales aplicadas.
    \item \textbf{Pruebas Automatizadas Unitarias e Integradas:} Ejecución determinista de casos de prueba en canalizaciones CI/CD antes de permitir el despliegue del prototipo.
    \item \textbf{Análisis Estático de Seguridad (SAST):} Escaneo automático de repositorios y contenedores para verificar la ausencia de credenciales expuestas y fallos de sintaxis.
    \item \textbf{Punto de Aprobación Humana (\textit{Human-in-the-Loop}):} Firma digital y revisión obligatoria por un investigador sénior antes de la puesta en producción.
    \item \textbf{Defensa Oral y Justificación Arquitectónica:} Instancia de evaluación donde los creadores del prototipo deben justificar las decisiones de diseño y demostrar la comprensión completa del sistema generado.
\end{enumerate}

\subsection{Plan Maestro e Implementación en Laboratorios de I+D}

La adopción institucional del prototipado acelerado se organiza en un Plan Maestro de cuatro fases proyectado a 12 meses:

\begin{itemize}
    \item \textbf{Fase 1: Gobernanza y Capacitación (Meses 1-2):} Formulación de la política institucional de uso de IA y bajo código; instalación de herramientas de análisis estático (SAST) y capacitación de investigadores en modelado UML y diseño de prompts.
    \item \textbf{Fase 2: Piloto de Automatización Documental (Meses 3-6):} Despliegue de flujos deterministas L0-L2 mediante Power Platform para simplificar la gestión de protocolos y aprobaciones de laboratorio, midiendo reducciones en el tiempo de ciclo.
    \item \textbf{Fase 3: Escalado de Microservicios y Agentes (Meses 7-10):} Adopción de FSMicroGenerator para proyectos de backend distribuido e integración de agentes de IA L3-L4 con supervisión humana estricta.
    \item \textbf{Fase 4: Auditoría, Medición y Consolidación (Meses 11-12):} Evaluación de impacto mediante indicadores KPI (cobertura funcional, deuda técnica, costo operativo) y consolidación de un catálogo de componentes reutilizables para la universidad.
\end{itemize}

% ==========================================
% NOTAS BIBLIOGRÁFICAS - CAPÍTULO 5
% ==========================================
\newpage
\markboth{CAPÍTULO 5\quad NOTAS BIBLIOGRÁFICAS E HISTÓRICAS}{}
\vspace*{0.8cm}
\noindent{\Huge \textbf{Notas bibliográficas\\[0.2cm] e históricas}}\\[1cm]

\textbf{Antecedentes:}\\[0.2cm]
El desarrollo del prototipado acelerado se fundamenta en la evolución desde las metodologías tradicionales (\textit{Waterfall}/Ágiles) hacia la Ingeniería Dirigida por Modelos (MDE/MDD). Históricamente, el desarrollo manual ha sufrido una alta sobrecarga sintáctica, donde hasta el 80\% del tiempo se destina a depuración y retrabajo evitable. La convergencia entre plataformas de bajo código e Inteligencia Artificial Generativa redefine este escenario mediante la \textbf{reasignación de tareas} (\textit{task reallocation}), desplazando la carga cognitiva humana desde la codificación repetitiva hacia la arquitectura, la verificación rigurosa y el cumplimiento normativo bajo estándares como ISO 15489 e ISO 27001.\\[0.5cm]

\textbf{Contribuciones Relevantes:}\\[0.2cm]
\begin{itemize}
    \item \textbf{Marian, Neferu y Mitrea (2026):} Aportan la evidencia del bloque determinista Low-Code en Power Platform, demostrando cuantitativamente una reducción del 65.6\% en los tiempos de aprobación documental y la optimización del TCO mediante listas estructuradas.
    \item \textbf{Khalfaoui, Khalfaoui y Azmani (2025):} Sustentan el pilar de generación automatizada de microservicios desde diagramas UML (FSMicroGenerator), logrando un 100\% de cobertura funcional en 3 horas sobre un stack MEAN.
    \item \textbf{Jauhari et al. (2026):} Proporcionan la Revisión Sistemática (PRISMA) sobre \textit{Vibe Coding}, evidenciando una aceleración del MVP del 40\% al 60\%, identificando el "cuello de botella de verificación" y justificando el marco VIBE4M.
    \item \textbf{Memari y Rudolph (2026):} Establecen la \textbf{Escalera de Capacidades} (L0 a L4) y la \textbf{Capa de Aseguramiento} (\textit{Assurance Layer}), definiendo protocolos de auditoría y validación \textit{Human-in-the-Loop}.
    \item \textbf{Sido y Emon (2024):} Suministran la metodología para el análisis de actores mediante la Matriz de Poder-Interés.
\end{itemize}
\vspace{0.5cm}

\textbf{Lecturas Sugeridas:}\\[0.2cm]
\begin{enumerate}
    \item \textbf{Brambilla, Cabot \& Wimmer (2017):} \textit{Model-Driven Software Engineering in Practice}. Base teórica formal sobre la transformación de modelos UML a código ejecutable.
    \item \textbf{ISO/IEC 15489-1:2016:} \textit{Information and documentation — Records management}. Estándar normativo para auditoría e inmutabilidad documental.
    \item \textbf{Maes (2025):} \textit{The gotchas of AI coding and vibe coding}. Guía práctica sobre el marco VIBE4M.
    \item \textbf{Zhao et al. (2025):} \textit{Is vibe coding safe? Benchmarking vulnerability of agent-generated code}. Estudio cuantitativo sobre vulnerabilidades de seguridad que justifica el escaneo SAST.
\end{enumerate}

\textbf{Referencias Completas}
\small
\begin{itemize}
    \item Brambilla, M., Cabot, J., \& Wimmer, M. (2017). \textit{Model-Driven Software Engineering in Practice} (2nd ed.). Morgan \& Claypool.
    \item ISO 15489-1:2016. (2016). \textit{Information and documentation — Records management}. ISO.
    \item Jauhari, A. F., et al. (2026). A systematic analysis of AI-assisted vibe coding in software development. \textit{IJST}, 5(2), 261–273.
    \item Khalfaoui, S., Khalfaoui, H., \& Azmani, A. (2025). Microservices-driven automation in full-stack development with FSMicroGenerator. \textit{IEEE Access}, 13, 110804–110825.
    \item Maes, S. H. (2025). \textit{The gotchas of AI coding and vibe coding: It's all about support and maintenance}. Zenodo.
    \item Marian, C. V., Neferu, M., \& Mitrea, D. A. (2026). Design and evaluation of a low-code/no-code document management system. \textit{Information}, 17(1), 46.
    \item Memari, M., \& Rudolph, G. (2026). The Capability Ladder: A curriculum-modernization framework for workforce readiness. \textit{arXiv}:2608.07779.
    \item Sido, N., \& Emon, E. A. (2024). \textit{Low/No Code Development and Generative AI} [Thesis Report, Aalborg University].
    \item Zhao, S., et al. (2025). Is vibe coding safe? Benchmarking vulnerability of agent-generated code. \textit{arXiv}:2512.03262.
\end{itemize}
\normalsize

\vspace{1cm}


% --- RESUMEN CAP 5 ---
\newpage
\markboth{CAPÍTULO 5\quad RESUMEN}{}
\vspace*{0.8cm}
{\Huge \textbf{Resumen del capítulo 5}}
\vspace{1cm}

\begin{resumenlist}
    \item \textbf{Orquestación Híbrida:} Integración del bajo código determinista, generadores UML y Vibe Coding en capas funcionales.
    \item \textbf{Selección y TCO:} Uso de matrices multidimensionales de selección para optimizar recursos en instituciones académicas.
    \item \textbf{Capa de Aseguramiento:} Validación \textit{Human-in-the-Loop} y análisis estático SAST para mitigar deuda técnica.
    \item \textbf{Plan Maestro:} Implementación sostenible en 4 fases para laboratorios de investigación y educación superior.
\end{resumenlist}

% --- CONCLUSIONES CAP 5 ---
\newpage
\markboth{CAPÍTULO 5\quad CONCLUSIONES}{}
\vspace*{0.8cm}
{\Huge \textbf{Conclusiones del capítulo 5}}
\vspace{1cm}

\subsubsection*{Conclusión principal}
La aceleración del prototipado científico mediante herramientas No-Code/Low-Code e Inteligencia Artificial Generativa no depende de la sustitución del código tradicional por una única tecnología, sino de la capacidad institucional para orquestar de manera híbrida soluciones deterministas y probabilísticas bajo un marco riguroso de gobernanza y aseguramiento de calidad.

\subsubsection*{Aportes y límites}
\textbf{Aportes:} Formulación de una matriz multidimensional de decisión, un algoritmo lógico de selección y una ruta de implementación en 4 fases adaptada a laboratorios de investigación. \textbf{Límites:} El éxito del modelo requiere un cambio cultural en los equipos de I+D y la disposición a invertir esfuerzo continuo en la verificación y auditoría de los sistemas generados.

\subsubsection*{Preguntas abiertas}
¿Cómo evolucionarán los estándares internacionales de auditoría de software a medida que los equipos de agentes autónomos de IA (Nivel L4) asuman la autogeneración y refactorización continua de sistemas en tiempo real?

% ==========================================
% REFERENCIAS BIBLIOGRÁFICAS
% ==========================================
\newpage
\markboth{REFERENCIAS BIBLIOGRÁFICAS}{REFERENCIAS BIBLIOGRÁFICAS} 
\phantomsection
\addcontentsline{toc}{section}{Referencias Bibliográficas}
\markboth{REFERENCIAS BIBLIOGRÁFICAS}{REFERENCIAS BIBLIOGRÁFICAS}
{\Huge \textbf{Referencias Bibliográficas}}\\[0.8cm]
\begin{itemize}[leftmargin=1.5em, label={}, itemsep=0.35ex, parsep=0pt]
\small
    \item Acemoğlu, D., Autor, D., Hazell, J., \& Restrepo, P. (2022). Artificial intelligence and jobs. \textit{Journal of Labor Economics}, 40(S1), S293–S340.
    \item Borgonovi, F., Calvino, F., Criscuolo, C., et al. (2023). Emerging trends in AI skill demand across 14 OECD countries. \textit{OECD Papers}.
    \item Brambilla, M., Cabot, J., \& Wimmer, M. (2017). \textit{Model-Driven Software Engineering in Practice} (2nd ed.).
    \item Evans, E. (2003). \textit{Domain-Driven Design: Tackling Complexity in Software}.
    \item Fowler, M. (2014). \textit{Microservices Architecture}.
    \item Fu, X., Bultan, T., \& Su, J. (2002). Formal verification of e-services and workflows.
    \item ISO 15489-1:2016. (2016). Information and documentation — Records management.
    \item ISO/IEC 27001:2022. (2022). Information Security Management Systems.
    \item Maes, S. H. (2025). The gotchas of AI coding and vibe coding: Support and maintenance.
    \item Manca, F. (2023). Six questions about the demand for artificial intelligence skills in labour markets. \textit{OECD Working Papers}.
    \item Martin, J. (1991). \textit{Rapid Application Development}. Macmillan Publishing Company.
    \item Meske, C., Hermanns, T., Weiden, E., Loser, K., \& Berger, T. (2025). Vibe coding as a reconfiguration of intent mediation in software development. \textit{IEEE Access}.
    \item Newman, S. (2015). \textit{Building Microservices: Designing Fine-Grained Systems}. O'Reilly Media.
    \item Ouyang, L., Wu, J., Jiang, X., et al. (2022). Training language models to follow instructions with human feedback. \textit{NeurIPS}, 35, 27730–27744.
    \item Prebreza, B., Gotseva, D., \& Nakov, P. (2021). A study of documents management system based on web, case study: University. \textit{IEEE TELECOM}, 85–89.
    \item Royce, W. W. (1970). Managing the development of large software systems: Concepts and techniques. \textit{IEEE WESCON}.
    \item Sartori, C. C. (2025). Architectures of error: A philosophical inquiry into AI and human code generation.
    \item Shi, Z., Dong, J., \& Gan, Y. (2025). Democratizing digital transformation: A multisector study of low-code adoption patterns. \textit{Applied Sciences}.
    \item van Hurne, M. H. K. (2025). A technical debt-aware prompting framework for sustainable vibe coding.
    \item Vaswani, A., Shazeer, N., Parmar, N., et al. (2017). Attention is all you need. \textit{NeurIPS}, 30, 5998–6008.
    \item Zhao, S., Wang, D., Zhang, K., et al. (2025). Benchmarking quality of agent-generated code in real-world tasks.
\end{itemize}

% ==========================================
% SOBRE ESTA OBRA
% ==========================================
\newpage
\markboth{SOBRE ESTA OBRA}{SOBRE ESTA OBRA} 
\phantomsection
\addcontentsline{toc}{section}{Sobre Esta Obra}
{\Huge \textbf{Sobre Esta Obra}}\\[1cm]
\textbf{Sinopsis:} \textit{Creación de Prototipos Rápidos: Guía de Herramientas No-Code/Low-Code para Validar Ideas Científicas} constituye una obra monográfica y un manual metodológico de referencia avanzada sobre la transformación digital de la investigación tecnológica y la ingeniería de software. Diseñada en estricto cumplimiento con la plantilla académica formal de la Universidad Nacional Mayor de San Marcos (UNMSM), esta obra examina la convergencia estratégica entre la abstracción visual de las plataformas Low-Code/No-Code deterministas, la generación automática de arquitecturas full-stack de microservicios en el stack MEAN a partir de diagramas de clases UML (mediante el marco FSMicroGenerator), la gobernanza documental bajo normas internacionales ISO 15489 e ISO 27001 en entornos Power Platform, la interacción conversacional en lenguaje natural (\textit{Vibe Coding}) y el control del código sintético impulsado por modelos de lenguaje de gran escala (LLMs).\\[0.5cm]
A través de cinco capítulos estructurados en profundidad, la obra introduce marcos de gobernanza de vanguardia como la Escalera de Capacidades (\textit{Capability Ladder}), la Capa Transversal de Aseguramiento (\textit{Assurance Layer}) y el protocolo VIBE4M. La obra está respaldada por la síntesis rigurosa y la evidencia cuantitativa de cinco investigaciones científicas internacionales fundamentales publicadas entre 2023 y 2026. Este volumen ofrece a docentes, investigadores, ingenieros de software y estudiantes universitarios un mapa de ruta replicable y metodológicamente validado para liderar el prototipado acelerado de software con la máxima velocidad, bajo costo infraestructural, control de calidad y rigurosidad académica.

\end{document}